\subsection{Associated RLA}%\label{assoc_rla}
On one hand, a RLA have associated a distribution given by the image of the anchor map. On the other hand, the image of the anchor map of a CLA is a complex distribution of $T_{\C}M$. To~bypass this issue, we associate two real distributions:
\begin{equation}\label{real_distributions}
\Delta=\re \bigl(\rho(A)\cap \overline{\rho(A)}\bigr)\qquad\text{and} \qquad D=\re \bigl(\rho(A) + \overline{\rho(A)}\bigr).
\end{equation}
Note that $D$ is a smooth distribution, while $\Delta$ is not smooth in general, see \cite[Example 6.1]{aguero2022complex}.

Given a CLA $(A, [\cdot,\cdot], \rho)$, decompose $\rho=\rho_1+{\rm i}\rho_2$. Contrary to what happens with CAVBs, neither $\mathcal{A}_1=(A_{\R}, [\cdot,\cdot], \rho_1)$ nor $\mathcal{A}_2=(A_{\R}, [\cdot,\cdot],\rho_2)$ necessarily define RLAs, since $\rho_1$ and $\rho_2$ do not preserve the bracket (and so they do not satisfy the Leibniz identity).


Actually, the natural candidate for $\mathcal{A}_2$ is $(A_{\R}, {\rm i}[\cdot,\cdot], -\rho_2)$, here we write $-\rho_2$ in order to avoid~${\rm i}\rho_2$. The failure of the Leibniz identity for $\mathcal{A}_1$ and $\mathcal{A}_2$ is controlled in the following way:
  \begin{gather}\label{fail_leib_re}
  [a, fb]_1-(f[a,b]_1+(\rho_1(a)f)b)={\rm i}(\rho_2(a)f)b,
  \\ \label{fail_leib_im}
  [a, fb]_2-(f[a,b]_2-(\rho_2(a)f)b)={\rm i}(\rho_1(a)f)b,
  \end{gather}
  where $[\cdot,\cdot]_1=[\cdot,\cdot]$ and $[\cdot,\cdot]_2={\rm i}[\cdot,\cdot]$.

An immediate consequence of equations \eqref{fail_leib_re} and \eqref{fail_leib_im} is the following.
\begin{Proposition}
Let $(A, [\cdot,\cdot], \rho)$ be a CLA. If either $\mathcal{A}_1=(A_{\R}, [\cdot,\cdot], \rho_1)$ or $\mathcal{A}_2=(A_{\R}, {\rm i}[\cdot,\cdot], -\rho_2)$ are RLAs, then $\rho=0$.
\end{Proposition}

In general, $\mathcal{A}_1$ and $\mathcal{A}_2$ satisfy the Jacobi identity but they are not RLAs, since there is an ``error term'' in the Leibniz identity controlled by the anchor of the other algebroid.

We avoid these error terms by working with $\ker\rho_1$ and $\ker\rho_2$. Also note that the multiplication by ${\rm i}$ is a kind of ``isomorphism'' from $\mathcal{A}_1$ to $\mathcal{A}_2$.

In the same way as with CAVBs, we consider the distribution
\[
A^{\real}=\ker \rho_2=\rho^{-1}(\Delta)
\]
and the real rank of a CLA as
\[\operatorname{\text{real-rank}} A\colon \ m\in M\mapsto \dim_{\R} A^{\real}|_m.\]
We also have the distribution given by $\ker\rho_1$. However, note that
\[\ker\rho_1=j_A(A^{\real})=\rho^{-1}({\rm i}\Delta),\]
 where $j_A$ is the multiplication by ${\rm i}$ in the fibers. Moreover,
\begin{equation}\label{id_1}
[A^{\real}, A^{\real}]\subseteq A^{\real},\qquad[A^{\real}, j_A(A^{\real})]\subseteq j_A(A^{\real})\qquad\text{and}\qquad [j_A(A^{\real}), j_A(A^{\real})]\subseteq A^{\real}.
\end{equation}
 Consider
 \[\rho^{\real}=\rho_1|_{A^{\real}}\qquad \text{and}\qquad \rho^{{\rm im}}=-\rho_2|_{j_A(A^{\real})}.\]

As a consequence, we get the following.
 \begin{Proposition}
 If $(A, [\cdot, \cdot], \rho)$ is a CLA with constant real rank, then $(A^{\real}, [\cdot, \cdot]_{|_{A^{\real}}}, \rho^{\real})$ and
 \smash{$(j_A(A^{\real}), j_A[\cdot, \cdot]_{|_{j_A(A^{\real})}}, \rho^{{\rm im}})$} are isomorphic RLAs, with the isomorphism given by the restric\-tion~of~$j_A$. Moreover, \smash{$\rho^{\real}(A^{\real})=\rho^{{\rm im}}(j_A(A^{\real}))=\Delta$}.
 \end{Proposition}

 {\samepage \begin{Definition}
 When $A^{\real}$ is a RLA, we shall call it as the {\em Lie algebroid of real elements}.
 \end{Definition}

 One of the key feature of the Lie algebroid of real elements is that it allows us to associate a~foliation to any CLA with constant real index.
 \begin{Corollary}
 If a CLA $(A, [\cdot, \cdot], \rho)$ has constant real index, then the set of real elements of the image of $\rho$, $\re\rho(A)$ coincides with the image of $\rho(A^{\real})$. Consequently, $\re\rho(A)$ is an integrable distribution.
 \end{Corollary}

 }

 A RLA is said to be {\em regular} when the image of the anchor map is a subbundle of $TM$. CLAs have three associated distributions: the image of the anchor map, $\Delta$ and $D$, see equation \eqref{real_distributions}. Therefore, a regularity condition for CLAs should take these two additional distributions into consideration.
 \begin{Definition}
 A CLA $(A,[\cdot,\cdot], \rho)$ is {\em regular} if $\rho(A)$ is a regular distribution and we say that it is {\em strongly regular} if $D$ and $\Delta$ are regular distributions.
\end{Definition}
 Note that strongly regular CLAs are regular CLAs with constant real rank, and vice versa.
\begin{Remark}
 Another natural invariant associated to a CLA is the {\em order}, which is a $\mathbb{N}$-function defined as follows:
 \[\order A\colon \ m\in M\mapsto \cork_{\R} D|_{m}.\]
The order was previously introduced in \cite{aguero2022complex} for complex Dirac structures and in \cite{tennyson2021exceptional}, under the name ``class'', for exceptional complex structures.
By equation \eqref{sum_identity}, we have that
\begin{equation}\label{orderrr}
\order A|_m=\rerk A|_m+\dim M-\rk A_{\R},
\end{equation}
and so, constant order is equivalent to constant real rank.
\end{Remark}

In general, the distribution $D$ is not integrable, even for strongly regular CLAs. Non-Levi-flat CR structures provide examples of strongly regular CLAs where the distribution $D$ is not integrable.
\begin{Definition}
Let $(A, [\cdot,\cdot], \rho)$ be a CLA (or an RLA) and let $m\in M$ be a point. The {\em isotropy Lie algebra of $A$ at $m$} is the complex (or real) Lie algebra
\[\g_m (A)=\ker\rho_m.\]
\end{Definition}
When $A$ is regular, we obtain a bundle of complex (or real) Lie algebras given by $\g(A)=\ker \rho$.
 Note that
 \begin{align*}
  \g(A)_{\R}= A^{\real}\cap j_A(A^{\real}) \qquad \text{and}\qquad
\g(A^{\real})=\g(j_A(A^{\real}))=\g(A)_{\R}.
   \end{align*}
